\documentclass[11pt,a4paper]{exam}
\usepackage[utf8]{inputenc}
\usepackage{amsmath, amsthm, amssymb, amsfonts} %ams packages
\usepackage{graphicx, color} %grphics
\usepackage{enumitem} %personalized lists 
\usepackage[top=3.5cm, bottom=2cm, left=2cm, right=2cm, headsep = 3pt, footskip=2pt]{geometry} %pagelayout
\usepackage{multirow,multicol} %to use multirow on tables
\usepackage{tabularx}
\usepackage{array} %bmatrix and similar stuff
% \usepackage{hyperref}
\usepackage[slovene]{babel}

\usepackage{multicol}
\usepackage{scalefnt}

%%%%%%%%%%%%%%%%%%%%%%%%%%%%%%%%%%%%%%%%%%%%%%%%%%%%%%%%%%%%%%%%%%%%%
%personalized commands to make latex look better (optional)
%%%%%%%%%%%%%%%%%%%%%%%%%%%%%%%%%%%%%%%%%%%%%%%%%%%%%%%%%%%%%%%%%%%%%%%%%%%
\renewcommand{\leq}{\leqslant} %menores y mayores iguales decentes
\renewcommand{\geq}{\geqslant}
\renewcommand{\subset}{\subseteq} 
\renewcommand{\supset}{\supseteq} 
\renewcommand{\subsetneq}{\varsubsetneq}
\renewcommand{\{}{\lbrace}
\renewcommand{\}}{\rbrace}
\newcommand{\sm}{\setminus} %setminus corto

%\renewcommand{\bar}{\overline}
%\renewcommand{\hat}{\widehat}

%%%%%%%%%%%%%%%%%%%%%%%%%%%%%%%%%%%%%%%%%%%%%%%%%%%%%%%%%%%%%%%%%%%%%}

%%%%%%%%%%%%%%%%%%%%%%%%%%%%%%%%%%%%%%%%%%%%%%%%%%%%%%%%%%%%%%%%%%%%%%%%%%%
%exam class commands
%%%%%%%%%%%%%%%%%%%%%%%%%%%%%%%%%%%%%%%%%%%%%%%%%%%%%%%%%%%%%%%%%%%%%%%%%%%
\pagestyle{head}
% \chead{ Matematična analiza 2024/2025 \\ Univerza v Ljublani, Pedagoška fakulteta \\ Vaje 1 \\17. februar 2025} %self-explanatory
% \runningheader{}{ Abstraktna algebra 2024/2025 \\ Univerza v Ljublani, Pedagoška fakulteta \\2. Kolokvij \\20. januar 2025}{} %self-explanatory
% \runningheader{}{}{}
% \footer{}{}{}
% \firstpagefooter{MATH 1190 3.0 M  }{Page \thepage\ of \numpages}{Copyright 2020 \textcopyright Antonio Montero}
% \firstpagefooter{}{}{}
% \runningfooter{MATH 1190 3.0 M }{Page \thepage\ of \numpages}{Copyright 2020 \textcopyright Antonio Montero}
% \firstpagefootrule 
% \runningfootrule

%%Solutions
\unframedsolutions
% \printanswers
\SolutionEmphasis{\itshape}


%Personalized commands
\pointpoints{točka}{točk}
\bonuspointpoints{točka}{dodatnih točk}
% \gradetable[h][questions]
\addpoints
\hqword{Naloga}
% \vpgword{text}
\hpword{Točke}
% \hsword{}
% \hbpword{}
\htword{Skupaj}

% \pointsinrightmargin
% \extrawidth{-2cm}
\marginpointname{ \points}


\chqword{Naloga}
\chpgword{Stran}
\chpword{točke}
% \chbpword{Puntos extra}
% \chsword{Puntos obtenidos
\chtword{Skupaj}
%%%%%%%%%%%%%%%%%%%%%%%%%%%%%%%%%%%%%%%%%%%%%%%%%%%%%%%%%%%%%%%%%%%%%

% \graphicspath{{../img/}} %Uso: \graphicspath{{RelativePath}}

%%%%%%%%%%%%%%%%%%%%%%%%%%%%%%%%%%%%%%%%%%%%%%%%%%%%%%%%%%%%%%%%%%%%%
%Personalized commands

%Personalized commands

\usepackage{tabularx}
\newcolumntype{C}{>{\centering\arraybackslash}X}%
\newcolumntype{R}{>{\raggedleft\arraybackslash}X}%
\newcolumntype{L}{>{\raggedright\arraybackslash}X}%

\newcolumntype{\cc}[1]{>{\hsize=#1\hsize\raggedright\arraybackslash}X}%
\newcolumntype{\rr}[1]{>{\hsize=#1\hsize\raggedleft\arraybackslash}X}%
\newcolumntype{\ll}[1]{>{\hsize=#1\hsize\centering\arraybackslash}X}%

\newcommand{\TT}{\mathbf{T}}
\newcommand{\FF}{\mathbf{F}}

\DeclareMathOperator{\sh}{sh}
\DeclareMathOperator{\ch}{ch}
\let\th\relax
\DeclareMathOperator{\th}{th}
\DeclareMathOperator{\arsh}{arsh}
\DeclareMathOperator{\arch}{arch}
\DeclareMathOperator{\arth}{arth}
\DeclareMathOperator{\arctg}{arctg}

%%%%%%%%%%%%%%%%%%%%%%%%%%%%%%%%%%%%%%%%%%%%%%%%%%%%%%%%%%%%%%%%%%%%%
		\newif\ifmarking
% 		\markingtrue %comment to hide marking suggestions

		\newenvironment{marking}
		{\hrulefill
		 
		\textbf{Marking suggestions:} 
		\ttfamily }
		{\hrulefill }

\allowdisplaybreaks

% \renewcommand{\thepartno}{\roman{partno}}
\renewcommand{\partlabel}{(\textit{\thepartno})}
%%%%%%%%%%%%%%%%%%%%%%%%%%%%%%%%%%%%%%%%%%%%%%%%%%%%%%%%%%%%%%%%%%%%%

\begin{document}
% \scalefont{1.3}
\chead{Matematična analiza 2024/2025 \\ Univerza v Ljubljani, Pedagoška fakulteta \\ 1. Kolokvij \\17. junij 2025}
% \footer{}{}{}

% \vspace{1 cm}

\pagestyle{head}
	

	
\begin{questions}
	\question Dana je funkcija $f: [0,2] \to \mathbb{R}$ s predpisom \[f(x) = \begin{cases}
    x^2 & \text{če je } x \in [0, 1) \\
    3-x & \text{če je } x \in [1, 2]
  \end{cases}\]
  \begin{parts}
    \part[12] Izračunajte zgornjo in spodnjo Riemannovo (oz. Darbouxjevo) vsoto funkcije $f$ na intervalu $[0, 2]$ glede na delitev $\mathcal{D}: 0 < \frac{1}{5} < \frac{3}{5} < 1 <  \frac{6}{5} < \frac{9}{5} < 2$.
    \part[13] Izračunajte zgornjo in spodnjo Riemannovo (oz. Darbouxjevo) vsoto funkcije $f$ na intervalu $[0, 2]$ glede na delitev $\mathcal{D}_{n}: 0 < \frac{1}{n} < \frac{2}{n} < \cdots <  \frac{2n-1}{n} < 2$. Nato po definiciji izračunajte določeni integral $\displaystyle \int_{0}^{2} f(x) \, dx$.
  \end{parts}
  \textbf{Nasvet:} Za poljubno naravno število $n$ velja $\sum_{i=0}^{n} i^{2} = \frac{n(n+1)(2n+1)}{6}$. Tega ni potrebno dokazovati, ampak lahko uporabite pri izračunu zgornje in spodnje Darbouxjeve vsote delitve $\mathcal{D}_{n}$.
	
  % \vspace{1 cm}



	\question Dano je zaporedje funkcij $f_n = x^n$ za vsak $n \in \mathbb{N}$.
  \begin{parts}
    \part[11] Dokažite, da zaporedje $\{f_n\}$ konvergira po točkah na intervalu $[0, 1]$ in določite limitno funkcijo.
    \part[14] Dokažite, da zaporedje $\left\{ f_n \right\}$ enakomerno konvergira na intervalu $\left[0, \frac{99}{100}\right]$.
    \bonuspart[10] Dokažite, da zaporedje $\{f_n\}$ ni enakomerno konvergentno na intervalu $[0, 1]$.
  \end{parts} 

  % \vspace{1cm}

\question Dani sta funkciji $f(x) = e^{x^{2}}$ in $g(x) = \cos(x)$.
  \begin{parts}
    \part[15] Izračunajte Taylorjeva razvoja funkcij $f$ in $g$ okoli točke $a=0$.
    \part[10] S pomočjo Taylorjevih razvojev funkcij $f$ in $g$ izračunajte naslednjo limito: \[\lim_{x \to 0} \frac{e^{x^{2}} -1}{1 - \cos(x)}.\]
  \end{parts}

	% \vspace{1cm}
	
  \question Dana je preslikava $d: \mathbb{R}^{2} \times \mathbb{R}^{2} \to \mathbb{R}$ s predpisom
  \[d((x_1, y_1), (x_2, y_2)) = \sqrt{|x_{1} - x_{2}|} + \sqrt{|y_{1} - y_{2}|}\]
  \begin{parts}
    \part[10] Dokažite, da je $d$ metrika na $\mathbb{R}^{2}$.
    \part[8] Skicirajte odprto kroglo $K_{d}((0, 0), 1)$ s polmerom $r = 1$ okoli točke $(0, 0)$ v metričnem prostoru $(\mathbb{R}^{2}, d)$.
    \part[7] Naj bo 
		%$(u,v) \in \mathbb{R}^{2}$ in 
		$T_{
		%(u,v)
		}: \mathbb{R}^{2} \to \mathbb{R}^{2}$ funkcija s predpisom 
		\[T_{
		%(u,v)
		}(x,y) = (x+
		1
		%u
		,y+
		1
		%v
		).\]
    Dokažite, da je $T_{
		%(1,2)
		}$ \emph{izometrija}, to pomeni, da ohranja razdalje, torej \[d(T_{
		%(1,2)
		}(x_1,y_1), T_{
		%(u,v)
		}(x_2,y_2)) = d((x_1,y_1), (x_2,y_2))\] za vse $(x_1,y_1), (x_2,y_2) \in \mathbb{R}^{2}$.
    %\bonuspart[7] Dokažite, da je $T_{(u,v)}$ izometrija za poljubno $(u,v) \in \mathbb{R}^{2}$.
    \bonuspart[10] Naj bo $d_{2}$ evklidska metrika na $\mathbb{R}^{2}$ in naj bo $K_{2}((0,0),r)$ odprta evklidska krogla s središčem $(0,0)$ in polmerom $r$, t.j.
		\[K_{2}((0,0),r) := \{(x,y) \in \mathbb{R}^{2} \mid \sqrt{x^2 + y^2} < r \}.\] 
    Dokažite, da obstaja taka konstanta $A > 0$, da je
    \[
		%\begin{aligned}
      K_{2}\left((0, 0), \frac{r}{A}\right) 
			%&
			\subseteq K_{d}((0, 0), r) 
			%&&\text{in} \\ K_{d}((0, 0), r) &\subseteq K_{2}((0, 0), r*B).
    %\end{aligned} 
		\]
    % \part[3] Dokažite da je podmnožica $A \subset \mathbb{R}^{2}$ odprta (oz. zaprta) v metričnem prostoru $(\mathbb{R}^{2}, d)$, če in samo če je odprta (oz. zaprta) v evklidskem metričnem prostoru $(\mathbb{R}^{2}, d_{2})$.
    \end{parts}
    % \textbf{Opombe:} Za rešitev katere koli točke te vaje lahko uporabite katero koli prejšnjo, tudi če je niste rešili.
\end{questions}

\clearpage

\chead{Matematična analiza 2024/2025 \\ Univerza v Ljubljani, Pedagoška fakulteta \\ 1. Izpit \\17. junij 2025}

\begin{questions}
\question Podana je funkcija \[f(x) = \frac{e^{x}}{e^{2x}+1}\]
\begin{parts}
  \part[13] Dokažite da je $f$ soda in s pomočjo odvoda skiciraj njen graf.
  \part[12] Izračunajte volumen vrtenine, ki nastane pri vrtenju grafa funkcije $f$ okoli osi $x$. 
\end{parts}

%\question[16] Iz kroga z radijem $1$ izrežemo izsek in preostanek zvijemo v stožec. Kolikšen naj bo kot izseka, da bo imel dobljeni stožec največji prostornino?

\question Izračunajte naslednje integrale
	\begin{parts}
		% \part[6] $\displaystyle \int \frac{x^2}{(1 + x^3)^2} \, dx$,
		\part[13] $\displaystyle \int \frac{3x + 5}{(x^2+1)(x-1)} \, dx$,
		\part[12] $\displaystyle  \int x^2 \ln(1 + x) \, dx$.

	\end{parts}


	\question Dano je zaporedje funkcij $f_n = x^n$ za vsak $n \in \mathbb{N}$.
  \begin{parts}
    \part[11] Dokažite, da zaporedje $\{f_n\}$ konvergira po točkah na intervalu $[0, 1]$ in določite limitno funkcijo.
    \part[14] Dokažite, da zaporedje $\left\{ f_n \right\}$ enakomerno konvergira na intervalu $\left[0, \frac{99}{100}\right]$.
    \bonuspart[10] Dokažite, da zaporedje $\{f_n\}$ ni enakomerno konvergentno na intervalu $[0, 1]$.
  \end{parts} 

  % \vspace{1cm}

%\question Dani sta funkciji $f(x) = e^{x^{2}}$ in $g(x) = \cos(x)$.
%  \begin{parts}
%    \part[10] Izračunajte Taylorjeve razvoja funkcij $f$ in $g$ okoli točke $a=0$.
%    \bonuspart[10] S pomočjo Taylorjevih razvojov funkcij $f$ in $g$, izračunajte naslednji limit: \[\lim_{x \to 0} \frac{e^{x^{2}} -1}{1 - \cos(x)}.\]
%  \end{parts}

\question Dana je preslikava $d: \mathbb{R}^{2} \times \mathbb{R}^{2} \to \mathbb{R}$ s predpisom
  \[d((x_1, y_1), (x_2, y_2)) = \sqrt{|x_{1} - x_{2}|} + \sqrt{|y_{1} - y_{2}|}\]
  \begin{parts}
    \part[10] Dokažite, da je $d$ metrika na $\mathbb{R}^{2}$.
    \part[8] Skicirajte odprto kroglo $K_{d}((0, 0), 1)$ s polmerom $r = 1$ okoli točke $(0, 0)$ v metričnem prostoru $(\mathbb{R}^{2}, d)$.
    \part[7] Naj bo 
		%$(u,v) \in \mathbb{R}^{2}$ in 
		$T_{
		%(u,v)
		}: \mathbb{R}^{2} \to \mathbb{R}^{2}$ funkcija s predpisom 
		\[T_{
		%(u,v)
		}(x,y) = (x+
		1
		%u
		,y+
		1
		%v
		).\]
    Dokažite, da je $T_{
		%(1,2)
		}$ \emph{izometrija}, to pomeni, da ohranja razdalje, torej \[d(T_{
		%(1,2)
		}(x_1,y_1), T_{
		%(u,v)
		}(x_2,y_2)) = d((x_1,y_1), (x_2,y_2))\] za vse $(x_1,y_1), (x_2,y_2) \in \mathbb{R}^{2}$.
    %\bonuspart[7] Dokažite, da je $T_{(u,v)}$ izometrija za poljubno $(u,v) \in \mathbb{R}^{2}$.
    \bonuspart[10] Naj bo $d_{2}$ evklidska metrika na $\mathbb{R}^{2}$ in naj bo $K_{2}((0,0),r)$ odprta evklidska krogla s središčem $(0,0)$ in polmerom $r$, t.j.
		\[K_{2}((0,0),r) := \{(x,y) \in \mathbb{R}^{2} \mid \sqrt{x^2 + y^2} < r \}.\] 
    Dokažite, da obstaja taka konstanta $A > 0$, da je
    \[
		%\begin{aligned}
      K_{2}\left((0, 0), \frac{r}{A}\right) 
			%&
			\subseteq K_{d}((0, 0), r) 
			%&&\text{in} \\ K_{d}((0, 0), r) &\subseteq K_{2}((0, 0), r*B).
    %\end{aligned} 
		\]
    % \part[3] Dokažite da je podmnožica $A \subset \mathbb{R}^{2}$ odprta (oz. zaprta) v metričnem prostoru $(\mathbb{R}^{2}, d)$, če in samo če je odprta (oz. zaprta) v evklidskem metričnem prostoru $(\mathbb{R}^{2}, d_{2})$.
    \end{parts}
    % \textbf{Opombe:} Za rešitev katere koli točke te vaje lahko uporabite katero koli prejšnjo, tudi če je niste rešili.
\end{questions}
\end{document}
